\chapter{蛋白质：折叠起来的分子机器}

\begin{kaichang}
找一枚生鸡蛋，打在透明玻璃杯里。蛋清是半透明的、稀稀的、会流动的液体。再把另一个鸡蛋煮熟剥开：蛋清变成了白色、结实、有弹性的一整块，把蛋黄稳稳包在中间。

这里有个值得停下来想一想的问题。水结成冰，加热还能化回水；巧克力熔化了，冷却又能凝固。可煮熟的蛋清，你把它放进冰箱、泡进水里、念任何咒语，它都再也不会变回那种透明的液体。加热到底对蛋清做了什么，让变化只朝一个方向走？

先记下你的猜测。这一章结束时，你不仅能回答这个问题，还会明白：你的头发、你的肌肉、你流的血、你伤口结的痂、你做豆腐的豆浆，全都是同一类分子变出来的戏法——蛋白质。
\end{kaichang}

\section{蛋清、头发和豆腐的共同点}

先看看生活中那些“蛋白质时刻”：

\begin{itemize}[itemsep=2pt]
\item 煮鸡蛋，蛋清从透明液体变成白色固体；
\item 理发店烫卷发，先涂药水、加热定型，直发能保持卷曲好几个月；
\item 豆浆里点一勺卤水（主要含氯化镁），整锅液体凝成豆腐；
\item 牛奶放久了变酸，结成絮状的小块；
\item 炖肉时，肉汤冷却后凝成肉冻，入口又化成鲜美的汤；
\item 划破手指，血很快凝住，挡住细菌。
\end{itemize}

这些现象看上去风马牛不相及，背后却站着同一类分子。蛋清里主要是卵清蛋白；头发和指甲主要是角蛋白；豆浆里是大豆蛋白；肉冻是煮散了的胶原蛋白；凝血靠的是纤维蛋白原。你的体重去掉水分之后，大约有一半是蛋白质——肌肉、皮肤、头发、骨骼里的支架、血液里运氧的血红蛋白，全在这本账上。

前两章我们认识了糖（第 46 章）和脂质（第 47 章）：糖像标准化的燃料块和建材，脂质像怕水的储能仓库和围墙材料。而蛋白质的角色完全不同：它们是\textbf{干活的机器}。糖可以被烧掉供能，脂质可以存起来备用，蛋白质却几乎每个都有具体职务——搬运、支撑、识别、防御、催化、收缩。一座城市不能只有燃料和砖，还得有起重机、传送带、门卫和扳手。蛋白质就是细胞这座城市的全套机械设备。

机器要有形状。起重机不能是一摊铁水，扳手不能是一团棉线。这一章的主线就是：蛋白质这台机器的形状从哪里来，又为什么会失去。

\section{二十种零件串成一根长链}

把镜头推到分子层。任何蛋白质，不管多复杂，拆到底都是一种零件的重复：\textbf{氨基酸}。

氨基酸这个名字已经泄露了它的结构（第 37 章讲过官能团的命名逻辑）：一端是氨基（\chem{{-}NH_2}，含氮，偏碱性），另一端是羧基（\chem{{-}COOH}，会给出氢离子，偏酸性）。这两头中间夹着一个碳原子，碳上还连着氢和一截\textbf{侧链}。所有氨基酸的“骨架”完全一样，差别全在侧链上——就像二十种插座相同的乐高零件，插口统一，外形各异。

人体用来造蛋白质的标准零件有 20 种。按侧链脾气，大致分几类：

\begin{itemize}[itemsep=2pt]
\item \textbf{怕水的}：侧链是碳氢链或碳氢环，跟油一个性格，见水就躲。亮氨酸、苯丙氨酸属于这类。
\item \textbf{亲水的}：侧链带羟基、氨基等极性基团，能和水拉氢键，愿意待在水环境里。丝氨酸、谷氨酰胺属于这类。
\item \textbf{带电的}：侧链在体液酸碱度下干脆带着正电或负电。赖氨酸带正电，谷氨酸带负电。
\item \textbf{有绝活的}：半胱氨酸侧链顶端有一个硫氢基（\chem{{-}SH}），两个半胱氨酸相遇时，两只硫可以拉起手来，形成一根牢固的共价桥——\textbf{二硫键}。记住它，你的卷发全靠它。
\end{itemize}

零件怎样串成链？靠第 38 章讲过的缩合反应：一个氨基酸的羧基和下一个氨基酸的氨基碰头，挤掉一个水分子，碳和氮之间形成一根共价键。这根键有个专门的名字：\textbf{肽键}。两个氨基酸缩合得二肽，三个得三肽，几十上百个连成\textbf{多肽链}——这就是蛋白质的本体。典型的蛋白质由几十到几千个氨基酸首尾相连，像一列挂着几百节车厢的火车，每节车厢是同一种底盘，载货各不相同。

链条反过来也能拆：肽键加水会水解（第 38 章），长链断开成游离的氨基酸。你的胃和小肠里天天在做这件事——这一点到“常见误解”那一节还要算总账。

顺便一提两件前情。第一，除甘氨酸外，氨基酸中间那个碳都是手性中心（第 42 章），而生命几乎只用左手型的那一半镜像——蛋白质是一台全部由左手零件装配的机器，所以它对别的分子的左右手极其挑剔。第二，蛋白质本质上是一种聚合物（第 43 章），和尼龙一样靠缩合反应成链；区别是尼龙只有一两种单体反复重复，蛋白质用 20 种单体排出独一无二的顺序。顺序，就是信息。

\section{长链怎样自己叠成机器}

一串软塌塌的链条，怎么会变成一台有固定形状的机器？答案是：它不是被谁捏出来的，而是\textbf{自己叠起来的}。折叠不需要图纸工人，只需要这条链、水，和第 27 章讲过的那条老规则——体系总是趋向自由能更低的安排。注意措辞：链条没有“想要”折叠的意图，它只是在亿万次随机扭动中，试出了能量上更划算的姿势，而能量划算的姿势出现得更频繁、停得更久。

推着链条叠起来的，是四股力量：

\begin{itemize}[itemsep=2pt]
\item \textbf{疏水效应}（主力）。怕水的侧链不愿暴露在水里。第 47 章讲过，疏水效应其实是水分子“嫌麻烦”：水宁愿让怕水的东西抱团缩进内部，自己好自由地拉氢键。于是链条像含羞草一样收拢，把怕水侧链藏进核心，亲水侧链留在外面——折叠的总方向盘。
\item \textbf{氢键}。链条骨架上的羰基氧和亚氨基氢互相拉氢键，把链的某些段落固定成规则的几何形状：绕成弹簧状的叫 \textbf{$\alpha$ 螺旋}，并排拉直成瓦片状的叫 \textbf{$\beta$ 折叠片}。这些重复出现的小结构叫\textbf{二级结构}。
\item \textbf{离子作用（盐桥）}。带正电的侧链和带负电的侧链隔空相吸，像两枚小磁铁，把链上相距很远的两处扣在一起。
\item \textbf{二硫键}。两个半胱氨酸的硫原子形成真正的共价键，是折叠结构上的“铆钉”，最耐高温耐拉扯。角蛋白里二硫键特别多，所以你的头发耐磨、蹄角坚硬。
\end{itemize}

折叠完成后，一根链盘成的整体形状叫\textbf{三级结构}——这才是单台机器的样子。有些机器还要几台拼起来才开工：几条折叠好的链（亚基）按固定方式组合，叫\textbf{四级结构}。血红蛋白就是四条链的合体。

于是蛋白质有了四个层次：一级是氨基酸的排列顺序，二级是局部的螺旋和折叠片，三级是整链的立体折叠，四级是多亚基的组装。这个层级表值得记住，但更要记住驱动它的那句话：\textbf{顺序决定折叠，折叠决定形状，形状决定功能}。

\begin{qiaoliang}
折叠有多不可思议？算一笔账。取一条 100 个氨基酸的链（在蛋白质里算短的），保守假设每个氨基酸只有 3 种可能的姿态。整条链的摆法就有
\[3^{100} \approx 5\times 10^{47}\ \text{种。}\]
假如链条每 $10^{-13}$ 秒换一种姿势（这已经快得离谱），试遍所有姿势要
\[5\times 10^{47}\times 10^{-13}\ \mathrm{s}=5\times 10^{34}\ \mathrm{s}\approx 1.6\times 10^{27}\ \text{年，}\]
而宇宙的年龄不过 $1.4\times 10^{10}$ 年。可是真实的蛋白质在毫秒到几秒之内就叠好了。结论只能是：折叠\textbf{不是穷举乱试}。链条每扭对一步，能量就降一截，正确的接触会稳定下来并引导下一步——像水流下山，不是随机走遍全地图才找到谷底，而是顺着坡度一路滑进山谷。这条“能量斜坡”是理解折叠的关键图景：一级序列不只是零件清单，它还编码了整条下坡路线。
\end{qiaoliang}

\section{给折叠画一张图纸}

符号最后登场。氨基酸的通用写法是把不变的骨架写出来，把可变的侧链记作 $\chem{R}$：

\[\chem{H_2N{-}CH(R){-}COOH}\]

肽键的形成，用反应式写出来就是缩合：

\[\chem{H_2N{-}CH(R_1){-}COOH} + \chem{H_2N{-}CH(R_2){-}COOH} \rightarrow \chem{H_2N{-}CH(R_1){-}CO{-}NH{-}CH(R_2){-}COOH} + \chem{H_2O}\]

中间新出现的 \chem{{-}CO{-}NH{-}} 就是肽键——它其实就是第 37 章见过的酰胺键，换了个蛋白质语境里的名字。肽键有个重要脾气：它偏平、不爱转动，像链条上一段段僵硬的小钢板，柔软只能来自钢板之间的连接处。折叠的几何学，一半是被这条“半刚性的链”限定的。

四个层次可以用一幅简图连起来：

\begin{center}
\begin{tikzpicture}[font=\small\sffamily]
\node[draw,rounded corners,align=center] (a) at (0,0) {氨基酸\\排成顺序};
\node[draw,rounded corners,align=center] (b) at (3.6,0) {螺旋与\\折叠片};
\node[draw,rounded corners,align=center] (c) at (7.2,0) {整链折叠\\成球};
\node[draw,rounded corners,align=center] (d) at (10.8,0) {几条链\\拼成机器};
\draw[->,thick] (a) -- node[above=1pt]{\scriptsize 氢键} (b);
\draw[->,thick] (b) -- node[above=1pt]{\scriptsize 侧链相互作用} (c);
\draw[->,thick] (c) -- node[above=1pt]{\scriptsize 组装} (d);
\node[anchor=north] at ([yshift=-4pt]a.south) {\scriptsize 一级};
\node[anchor=north] at ([yshift=-4pt]b.south) {\scriptsize 二级};
\node[anchor=north] at ([yshift=-4pt]c.south) {\scriptsize 三级};
\node[anchor=north] at ([yshift=-4pt]d.south) {\scriptsize 四级};
\end{tikzpicture}
\end{center}

这是人为的表示法：方框和箭头只是为了帮助思考，真实的折叠没有这么整齐的舞台，而是一锅所有过程同时进行的分子热运动。

\section{形状就是本领}

“结构决定功能”这句话在蛋白质身上体现得最淋漓尽致。看五台机器：

\textbf{血红蛋白}是快递员。四条链（四级结构）各抱一个血红素，血红素中心的铁原子（就是第 10 章开场提到的那个铁）轻轻扣住一个氧分子。在肺里氧多，它抓氧；在组织里氧少，它放氧。更妙的是，一条链抓到氧后形状微微一变，会“推”旁边三条链一把，让它们抓氧更利索——这叫协同效应，是四级结构才有的团队技能。

\textbf{胶原蛋白}是钢缆。它把三条链像麻绳一样拧成一股三股螺旋，关键诀窍是每隔两个氨基酸必须出现一个甘氨酸——最小、没有侧链负担的氨基酸，因为绳子的正中心只容得下它。皮肤、肌腱、骨骼的韧性都靠这种缆绳。老火靓汤里的胶冻，就是三股螺旋被长时间加热拆散后的产物。

\textbf{抗体}是门卫的对讲机。Y 形的分子，两个“手臂”顶端是能识别特定外来分子（抗原）的接口，几十亿种抗体各认各的目标；“尾巴”则负责呼叫免疫细胞来处理。你的身体能为几乎任何从没见过的入侵者现场定制抗体——这是第三册免疫章节的故事，这里只需记住：识别的精准度，来自折叠出来的形状之间严丝合缝的互补。

\textbf{肌动蛋白}是可拆装的脚手架。单个肌动蛋白是球形的，却能像串珠一样首尾相连，拼成长长的纤维，需要时又能拆开。肌肉收缩、细胞爬行、细胞分裂时一分为二，都靠这种“会组装的绳子”和它的搭档马达蛋白。

\textbf{离子通道}是带门禁的水管。它横插在细胞膜上，折叠出一个贯穿膜的孔道，孔道内壁的电荷和宽窄只允许特定的离子（比如钠离子或钾离子）通过，还能开能关。你的神经能传递信号，靠的就是这些通道的开关节奏——第 51 章讲细胞膜时会细说，这里先埋个位置。

最后要修正一个可能的错觉：折叠好的蛋白质不是一尊凝固的雕塑。它在热运动中不断轻微抖动、开合，像一台运转中的机器而不是陈列品。血红蛋白抓放氧时会变形，通道开关时会变形，酶结合底物时也会变形（第 49 章的“诱导契合”还会展开）。结构决定功能，但\textbf{会变化的结构}才决定活的功能。甚至有些蛋白质的功能区天生就不折叠，像一截灵活的绳段，专门用来同时勾住好几个伙伴——“好蛋白质必须叠成固定形状”这条规则，在这里自己就破了例。

\section{把机器拆散，它能自己装回去吗}

折叠“由序列自动完成”，是个很大胆的说法。凭什么相信它，而不是相信“细胞里有某种装配车间，把链捏成形状”？

\begin{zhengju}
20 世纪 50 年代，美国生物化学家安芬森（Christian Anfinsen）盯上了一种叫核糖核酸酶的小蛋白质：只有 124 个氨基酸，链内有 4 对二硫键，能把 RNA 剪碎，活性很容易测量。他的实验设计朴素得惊人，分两步。

第一步，拆。把核糖核酸酶泡在浓尿素溶液里（尿素会搅乱氢键和疏水环境），再加一种试剂把 4 对二硫键全部剪断。蛋白质彻底摊开成一根无序的软链，活性归零——机器被拆成了零件。

第二步，装。用透析的办法慢慢把尿素和试剂滤走，只留下蛋白质和空气。几小时后检测：蛋白质重新叠回了原来的形状，4 对二硫键准确配对，活性几乎完全恢复。

关键在数字。4 对二硫键，从 8 个半胱氨酸里挑配对，一共有 105 种配法，只有 1 种是对的。如果复性靠乱碰，正确率应该只有约 1\%——但实验里恢复的是几乎全部的活性。安芬森还做过对照：如果先让链在“错误环境”里随便配对二硫键，再纠正，活性恢复得又慢又差。这说明正确的折叠信息不在外界的模具里，也不在二硫键里，而是\textbf{写在氨基酸的顺序本身}：只要序列还在、环境正常，折叠自动朝正确的谷底走。这就是“一级结构决定高级结构”的实验根基，安芬森因此分享了 1972 年诺贝尔化学奖。

这个故事也有后续修正：后来在真正的细胞里发现，很多蛋白质折叠时需要别的蛋白质（伴侣蛋白）在旁边“扶一把”，防止中途粘连。所以严谨的说法是：序列编码了折叠的终点和路线，但拥挤的细胞环境里，这条路有时需要交警。科学结论不是被推翻，而是被放进了更真实的环境里。
\end{zhengju}

\section{叠错了会怎样}

折叠是可逆拆装的平衡游戏，那么把平衡推开，机器就散架了。这就是\textbf{变性}：加热、强酸、强碱、酒精、重金属离子、剧烈振荡，都能破坏维持折叠的弱相互作用（氢键、离子作用、疏水环境），让链条摊开。注意一个要点：\textbf{变性拆的是折叠，不是链条本身}——一级结构的肽键在蒸煮酸泡中大多完好，所以变性蛋白质并没有“分解成别的东西”，只是机器散了架。

\begin{shiyan}
\textbf{观察：让蛋清变性给你看}（约 20 分钟）

材料：1 枚鸡蛋、3 个透明小杯、白醋、热水（约 $60\,^{\circ}\mathrm{C}$，手摸着烫但不至于烫伤）、筷子、标签纸。

步骤：
\begin{enumerate}[itemsep=1pt]
\item 把蛋清和蛋黄分开，蛋清搅匀后均分到三个杯子，贴上标签：对照、加酸、加热。
\item “加酸”杯里倒入约三分之一杯白醋，搅几下，静置 10 分钟。
\item “加热”杯放进盛热水的大碗里隔水浴 10 分钟（不要让蛋清沸腾）。
\item “对照”杯什么都不做。
\end{enumerate}

观察点：三个杯子的透明度、流动性各有什么变化？把“加酸”杯和“加热”杯再静置半小时，变化会退回吗？为什么两个完全不同的处理，得到的现象却相似？

安全提示：全程只用食品级材料，热水由大人协助准备；实验后的蛋清\textbf{不要吃}（接触过醋和长时间室温放置）；生蛋清可能带沙门氏菌，做完认真洗手，别让它碰到熟食和餐具。
\end{shiyan}

变性能可逆吗？安芬森的实验说“能”，煮熟的鸡蛋说“不能”。区别在\textbf{聚集}：实验室里复性的是稀溶液里孤零零的分子，而熟蛋清里无数条摊开的链挤在一起，疏水表面互相缠成一团乱麻，还拉起新的相互作用——就像一万根耳机线倒进一个箱子。摊开可逆，缠死不可逆。所以“变性不可逆”不是定律，是条件：安芬森能复原他的酶，你复原不了你的鸡蛋，都对。

细胞自己也要面对叠错的问题。刚出生的肽链一边合成一边折叠，周围拥挤不堪，很容易中途粘连。细胞的对策是\textbf{伴侣蛋白}：一类专门“陪产”的蛋白质，有的像桶一样把新生链罩起来让它安静折叠，有的把叠错的链拆开重试。这个过程要消耗能量（由 ATP 提供，第 50 章细讲）——你看，连“正确地叠起来”这件看似自动的事，在生命的账本上也是要花钱的。

最怕的是叠错还赖着不走。有些错误折叠的蛋白质会聚成不溶的硬块，细胞拆不掉。阿尔茨海默病患者脑中的淀粉样斑块，就是一类叠错的蛋白质层层堆积的结果（具体因果仍在研究中，科学界还没有完整答案，这是实话）。更离奇的是朊病毒：一种蛋白质叠成错误形状后，居然能“劝说”正常形状的同类变成错误形状，像错误折叠的传染病——疯牛病就是这样传播的。它没有 DNA、没有细胞结构，只是一团形状错误的蛋白质，却能让错误自我复制。这个发现当年轰动科学界，因为它打破了“自我复制必须靠核酸”的铁律。

\begin{wuqu}
“吃胶原蛋白能补皮肤、补关节，吃什么补什么。”——这是商家最爱利用的直觉。真相在消化道上：胶原蛋白进了胃肠，会被拆成氨基酸和小肽（第 49 章讲的消化酶干的就是这件事），和你吃的鸡蛋、豆腐里的蛋白质殊途同归。身体再拿这些零件，按\textbf{自己的}图纸组装它需要的东西——不会因为原料姓“胶原”，就定向送到你的皮肤上。否则吃猪脑补脑、吃核桃补脑早就该灵了。胶原蛋白口服液不比其他优质蛋白更“补皮肤”，价格却往往贵得多。同理，“敷面膜让皮肤吸收胶原蛋白”也走不通：完整的蛋白质分子太大，穿不过皮肤屏障。想让身体合成胶原蛋白，真正管用的配料是普通优质蛋白、维生素 C（胶原合成的必需助手）、睡眠和防晒。
\end{wuqu}

\section{回到那枚鸡蛋}

现在可以清算开场的悬念了。

生蛋清里，卵清蛋白分子各自折叠成小球，球的表面是亲水侧链，所以它们安分地分散在水里，互不纠缠。小球比光的波长小得多，光线穿过去基本不拐弯——蛋清看起来透明、流动。

加热后，热运动越来越猛，弱相互作用撑不住了：氢键断开，折叠松开，藏在核心的疏水侧链暴露到水里。水“拒绝”接待它们，它们只好去抱同伴的大腿——一条链的疏水补丁贴上另一条链的疏水补丁。亿万条链这样互相勾连，结成一张贯穿整杯蛋清的三维大网，把水锁在网格里。液体就此变成固体。大网的线团尺寸赶上了光的波长，光线被撞得四处散射——蛋清看起来白了。这就是那个“白色”：不是添加了什么色素，而是结构尺度变化带来的光学效果，和牛奶是白色、浪花是白色一个道理。

为什么回不去？因为成的网不是单分子的折叠错误，而是亿万分子缠死的聚集态，加上加热还促成了一些新的牢固接触（包括二硫键的重新配对）。要让一切复原，得把所有链同时拆开、稀释、再各自慢慢折好——冰箱显然没这个本事。安芬森的试管里做得到，你的锅里做不到，差别只在浓度和纠缠。

还有一个更有意思的对照：生鸡蛋的蛋白质不仅是“能吃”，还是“活的材料”——受精的鸡蛋里，这些分子机器配合起来能造出一只小鸡。煮熟之后，氨基酸几乎一个没少（营养上熟鸡蛋甚至更好消化，因为摊开的链把肽键都亮给了消化酶），但那台“能造鸡的机器”永久停机了。烧掉的不是物质，是\textbf{组织起来的结构}。这正是第 45 章那句话的分子版：生命的本质不在原料清单，在组织方式。

\section{机器还要上电才会转}

蛋白质化学早就走出了课本。医院里，单克隆抗体药物是能精确打击癌细胞表面标志的“定制导弹”；糖尿病患者注射的胰岛素，是用工程菌批量生产的蛋白质（第 83 章会讲这条生产线怎样改写医学）。理发店里，烫发的全套流程就是二硫键化学：先用药水拆断角蛋白里的二硫键，让头发“散架”变软；卷出想要的形状；再用氧化剂让二硫键在新位置上重新铆合——发型是被共价键记住的。食品工业里，凝乳酶让牛奶凝成奶酪，木瓜蛋白酶让老肉变嫩，豆腐坊的卤水用镁离子中和蛋白质表面的电荷、拆掉它们的“互斥护栏”，逼它们凝聚。

至于这些机器里最神奇的一类——能让化学反应加速百万倍、却只催促不消费自身的\textbf{酶}——它们绝大多数也是蛋白质，折叠出的那个小口袋（活性中心）凭什么有这么大神通？它降低的到底是什么，又绝不能改变什么？这正是下一章的主题。我们刚刚修好了机器的外壳，接下来该点火启动了。

\begin{tiaozhan}
折叠态到底比摊开态“划算”多少？答案是：只划算一点点。一条典型蛋白质，折叠态的自由能只比摊开态低约 $20$ 到 $60\ \mathrm{kJ\,mol^{-1}}$——摊到上百个氨基酸头上，每个残基还不到半个氢键的分量。要知道，整条链上起作用的氢键、盐桥、疏水接触数以百计，推力和拉力都是天文数字，净结果却在刀刃上。这种“边缘稳定”不是设计缺陷，而是特性：差太多，环境一波动就散架；稳太多，机器就僵死、没法变形干活、报废时也没法拆解回收。蛋白质的工作点，恰好落在稳定与灵活的刀锋之间。如果你学了第三册的进化篇再回头看，会更能体会：这把刀锋不是谁调出来的，而是几十亿年试错筛选出的可行区间。数学上，折叠的方向由 $\Delta G=\Delta H-T\Delta S$ 裁定：焓变（成键、氢键，放能有利）与熵变（链变有序不利、水被解放有利）在此博弈——疏水效应之所以是主力，恰恰因为它是“水的熵增”这一头的筹码。
\end{tiaozhan}

\section*{练一练}

\begin{sikaoti}
\item 画两个氨基酸（侧链分别标 $\chem{R_1}$、$\chem{R_2}$）缩合成二肽的粒子图：标出氨基、羧基、脱去的水分子来自哪两个基团、新生成的肽键在哪里。再用同一幅图的倒放，说明胃里发生了什么。
\item 某蛋白质核心的一个怕水侧链被换成了带电侧链（其余不变）。预测：这个蛋白质折叠后会发生什么？它在细胞里可能引发什么后果？请从四股折叠力量逐一分析。
\item 参考本章的“能量斜坡”图景，画一幅能量示意图：横轴是折叠进度，纵轴是自由能，标出摊开态、折叠态，并解释为什么折叠不需要翻越大山。再想一想：酶催化的反应坐标图（第 27 章）与此图有何不同？
\item 烫发药水分“拆键”和“定型”两步。用二硫键的语言解释这两步各自做了什么；并预测：天生卷发的人，角蛋白里的二硫键分布与直发者可能有何差别？
\item 点豆腐用卤水（含镁离子），不用食盐水也能成，但用大量纯水冲稀豆浆却永远不会凝固。从蛋白质表面电荷和聚集的角度，给出一个解释。
\item 生蛋清透明、熟蛋清发白，都不涉及颜色物质的变化。再举出两个“结构尺度改变光学效果”的生活例子（提示：牛奶、浪花、云），并说明它们与熟蛋清发白的共同原理。
\end{sikaoti}
